\documentclass[11pt,a4paper]{article}
\input{../../sablony/pripravy/priprava.sty}
\nastavPripravu{1. ročník}{Kvadratické rovnice a nerovnice}{125}{Rovnice v podílovém tvaru}
\begin{document}
\blok{Co se změní, když je ve jmenovateli neznámá?}
Oproti běžné kvadratické rovnici musíme ještě před úpravami určit
\textbf{definiční podmínky}. Krácením nesmí do množiny řešení přibýt
zakázaná hodnota. Rovnice $\frac{A(x)}{B(x)}=0$ platí právě tehdy,
když $A(x)=0$ a současně $B(x)\ne0$.
\blok{Společně řešené příklady -- postupně těžší}
\begin{prikladReseniBox}[title={Příklad 1 -- přímé krácení}]
Řešte $\frac{x^2-9}{x-3}=5$.
\tcblower
\[\begin{aligned}x&\ne3\\\frac{(x-3)(x+3)}{x-3}&=5\\x+3&=5\\x&=2\\K&=\{2\}\end{aligned}\]
\end{prikladReseniBox}
\begin{prikladReseniBox}[title={Příklad 2 -- zakázaný kořen}]
Řešte $\frac{x^2-5x+6}{x-2}=0$.
\tcblower
\[\begin{aligned}x&\ne2\\\frac{(x-2)(x-3)}{x-2}&=0\\x-3&=0\\K&=\{3\}\end{aligned}\]
Nulový čitatel má také kořen $2$, ale ten je ve jmenovateli zakázaný.
\end{prikladReseniBox}
\begin{prikladReseniBox}[title={Příklad 3 -- nekonečně mnoho řešení}]
Řešte $\frac{x^2-4}{x+2}=x-2$.
\tcblower
\[\begin{aligned}x&\ne-2\\\frac{(x-2)(x+2)}{x+2}&=x-2\\x-2&=x-2\\K&=\mathbb R\setminus\{-2\}\end{aligned}\]
\end{prikladReseniBox}
\blok{Samostatné procvičování}
\begin{zadaniReseniBox}[title={Úloha 1}]
$\frac{x^2-16}{x-4}=x+4$.
\tcblower
\[\begin{aligned}x&\ne4\\x+4&=x+4\\K&=\mathbb R\setminus\{4\}\end{aligned}\]
\end{zadaniReseniBox}
\begin{zadaniReseniBox}[title={Úloha 2}]
$\frac{x^2-1}{x+1}=0$.
\tcblower
\[\begin{aligned}x&\ne-1\\\frac{(x-1)(x+1)}{x+1}&=0\\x&=1\\K&=\{1\}\end{aligned}\]
\end{zadaniReseniBox}
\blok{Domácí úkol}
\begin{zadaniReseniBox}[title={Domácí příprava}]
\begin{enumerate}[label=\alph*)]
\item $\frac{x^2-25}{x-5}=7$.
\item $\frac{x^2-16}{x+4}=0$.
\item $\frac{x^2-4}{x-2}=x+2$.
\item $\frac{x^2-4x+3}{x-1}=0$.
\item $\frac{x^2-x-6}{x+2}=0$.
\item $\frac{x^2-9}{x-3}=x+3$.
\end{enumerate}
\tcblower
\begin{enumerate}[label=\alph*)]
\item $x\ne5$, zkrátíme na $x+5=7$. $K=\{2\}$.
\item $x\ne-4$, zkrátíme na $x-4=0$. $K=\{4\}$.
\item $x\ne2$, identita pro všechna přípustná $x$; $K=\mathbb R\setminus\{2\}$.
\item $x\ne1$, zkrátíme na $x-3=0$. $K=\{3\}$.
\item $x\ne-2$, zkrátíme na $x-3=0$. $K=\{3\}$.
\item $x\ne3$, identita; $K=\mathbb R\setminus\{3\}$.
\end{enumerate}
\end{zadaniReseniBox}
\blok{Dobrovolné rozšíření -- podílové nerovnice}
Zde \emph{nemůžeme} násobit jmenovatelem bez znalosti jeho znaménka.
Kritické body z čitatele a jmenovatele vyznačíme do znaménkové tabulky.
\begin{zadaniReseniBox}[title={Ukázková nerovnice}]
$\frac{(x-2)(x+2)}{x-1}>0$.
\tcblower
Kritické body jsou $-2,1,2$. Znaménka na čtyřech intervalech:
\[\begin{array}{c|cccc}
x&(-\infty;-2)&(-2;1)&(1;2)&(2;+\infty)\\\hline
\frac{(x-2)(x+2)}{x-1}&-&+&-&+
\end{array}\]
\[K=(-2;1)\cup(2;+\infty).\]
\end{zadaniReseniBox}
\poznamka{Původní úlohy na podílové nerovnice zachovány níže jako samostatný \textbf{dobrovolný} blok; nejsou povinným obsahem nové 125. hodiny.}
\blok{Původní náročné úlohy -- jen pro zájemce}
\begin{zadaniReseniBox}[title={Domácí příprava}]
\begin{enumerate}[label=\alph*),itemsep=2pt]
\item $\dfrac{\left(x - 2\right) \cdot \left(x + 2\right)}{x - 1}>0$
\item $\dfrac{\left(x - 3\right) \cdot \left(x + 4\right)}{x - 2}\le0$
\item $\dfrac{\left(x - 5\right) \cdot \left(x + 1\right)}{x}\ge0$
\item $\dfrac{\left(x - 1\right) \cdot \left(x + 3\right)}{x - 4}<0$
\item $\dfrac{\left(x - 3\right) \cdot \left(x - 1\right)}{x + 1}\le0$
\item $\dfrac{\left(x - 5\right) \cdot \left(x + 2\right)}{x - 2}>0$
\end{enumerate}
\tcblower
\begin{enumerate}[label=\alph*),itemsep=2pt]
\item
\[\begin{aligned}
x&\ne 1\quad\text{(podmínka)}\\[2pt]
\text{Kritické body: }&-2,\ 1,\ 2\\[2pt]
K&=(-2;1)\cup(2;+\infty)
\end{aligned}\]
\item
\[\begin{aligned}
x&\ne 2\quad\text{(podmínka)}\\[2pt]
\text{Kritické body: }&-4,\ 2,\ 3\\[2pt]
K&=(-\infty;-4\rangle\cup(2;3\rangle
\end{aligned}\]
\item
\[\begin{aligned}
x&\ne 0\quad\text{(podmínka)}\\[2pt]
\text{Kritické body: }&-1,\ 0,\ 5\\[2pt]
K&=\langle-1;0)\cup\langle5;+\infty)
\end{aligned}\]
\item
\[\begin{aligned}
x&\ne 4\quad\text{(podmínka)}\\[2pt]
\text{Kritické body: }&-3,\ 1,\ 4\\[2pt]
K&=(-\infty;-3)\cup(1;4)
\end{aligned}\]
\item
\[\begin{aligned}
x&\ne -1\quad\text{(podmínka)}\\[2pt]
\text{Kritické body: }&-1,\ 1,\ 3\\[2pt]
K&=(-\infty;-1)\cup\langle1;3\rangle
\end{aligned}\]
\item
\[\begin{aligned}
x&\ne 2\quad\text{(podmínka)}\\[2pt]
\text{Kritické body: }&-2,\ 2,\ 5\\[2pt]
K&=(-2;2)\cup(5;+\infty)
\end{aligned}\]
\end{enumerate}
\end{zadaniReseniBox}
\end{document}
